% ============================================================================
% pilotaz-cennik-umowa.tex -- Pilotaż 12 miesięcy: cennik i wzór umowy
%
% Rodzaj dokumentu: oferta i wzór umowy dla gabinetów (B2B), a nie praca
% badawcza w rozumieniu docs/research/00-RESEARCH-DOCUMENTATION.md. Nie wprowadza
% liczb ani faktów spoza tego pliku. Wspólna preambuła house-preamble.tex daje
% tylko wygląd.
%
% \paperstatus jest tu użyty na inne zastosowanie niż w szablonie badawczym:
% oznacza status wersji dokumentu (wersja robocza, do uzupełnienia danych).
%
% Dokument zawiera pola w nawiasach [ ] do uzupełnienia przed podpisaniem.
% Build i publikacja PDF: .github/workflows/deploy-pages.yml.
% ============================================================================
\documentclass[11pt,a4paper]{article}
\newcommand{\paperstatus}{WERSJA ROBOCZA}
\newcommand{\headerleft}{Inference Polska --- Pilotaż: cennik i wzór umowy}
\newcommand{\pdftitleline}{Inference Polska: pilotaż 12 miesięcy, cennik i wzór umowy}
\newcommand{\pdfauthorline}{Konrad Cinkusz}
\def\houselang{polish}
% ============================================================================
% HOUSE-PREAMBLE.tex -- the estate's house LaTeX style, as a file you copy in
% and \input rather than a block you copy out.
%
% See 00-RESEARCH-DOCUMENTATION.md, "The house preamble is a file, not a
% block". Copy this file to <repo>/docs/papers/house-preamble.tex (or
% docs/research/papers/ for a study paper) and \input it from every document
% in the repository that wants the house look. One copy per repository; every
% document in that repository \inputs the same one.
%
% WHY A FILE. The estate's first attempt at this shipped a template whose
% preamble was meant to be copied into each new document. Copying is what
% drifts: `ab-ove`'s two overview editions carried none of the seven house
% markers -- not the colors, the section formatting, the header, the status
% marker, the path macro, the geometry or the link style -- and nothing said
% so, because a convention about copying correctly has nothing to fail. The
% Beamer theme, which was always distributed as a file, propagated
% byte-identically over the same period. Distribute the file.
%
% ----------------------------------------------------------------------------
% THE CONTRACT. The including document defines these BEFORE \input, and
% carries its own \documentclass and \begin{document} -- this file has
% neither, so that one preamble can serve documents with different classes.
%
%   \paperstatus    Right-hand header, and the title block's status line.
%                   A study paper: DRAFT or VERIFIED, mirroring the study's
%                   own header. Any other document: its own status, named in
%                   that document's header comment as repurposed.
%   \headerleft     Left-hand header. Usually "Project --- Document title".
%   \pdftitleline   PDF metadata title.
%   \pdfauthorline  PDF metadata author.
%
% Optional, both defaulted below if the document does not define them:
%
%   \houselang      babel language for this edition, e.g. \def\houselang{polish}.
%                   Defaults to english. Set it in the edition file, not here:
%                   it is the one thing that is genuinely per-edition.
%   \editionsuffix  This edition's diagram-file suffix, e.g.
%                   \def\editionsuffix{.pl}. Defaults to empty. \dgm reads it,
%                   so a figure is called by slug and resolves per edition --
%                   see "Figures" below.
%
% A minimal English document:
%
%   \documentclass[11pt,a4paper]{article}
%   \newcommand{\paperstatus}{DRAFT}
%   \newcommand{\headerleft}{Project --- Overview}
%   \newcommand{\pdftitleline}{Project: an overview}
%   \newcommand{\pdfauthorline}{Author Name}
%   % ============================================================================
% HOUSE-PREAMBLE.tex -- the estate's house LaTeX style, as a file you copy in
% and \input rather than a block you copy out.
%
% See 00-RESEARCH-DOCUMENTATION.md, "The house preamble is a file, not a
% block". Copy this file to <repo>/docs/papers/house-preamble.tex (or
% docs/research/papers/ for a study paper) and \input it from every document
% in the repository that wants the house look. One copy per repository; every
% document in that repository \inputs the same one.
%
% WHY A FILE. The estate's first attempt at this shipped a template whose
% preamble was meant to be copied into each new document. Copying is what
% drifts: `ab-ove`'s two overview editions carried none of the seven house
% markers -- not the colors, the section formatting, the header, the status
% marker, the path macro, the geometry or the link style -- and nothing said
% so, because a convention about copying correctly has nothing to fail. The
% Beamer theme, which was always distributed as a file, propagated
% byte-identically over the same period. Distribute the file.
%
% ----------------------------------------------------------------------------
% THE CONTRACT. The including document defines these BEFORE \input, and
% carries its own \documentclass and \begin{document} -- this file has
% neither, so that one preamble can serve documents with different classes.
%
%   \paperstatus    Right-hand header, and the title block's status line.
%                   A study paper: DRAFT or VERIFIED, mirroring the study's
%                   own header. Any other document: its own status, named in
%                   that document's header comment as repurposed.
%   \headerleft     Left-hand header. Usually "Project --- Document title".
%   \pdftitleline   PDF metadata title.
%   \pdfauthorline  PDF metadata author.
%
% Optional, both defaulted below if the document does not define them:
%
%   \houselang      babel language for this edition, e.g. \def\houselang{polish}.
%                   Defaults to english. Set it in the edition file, not here:
%                   it is the one thing that is genuinely per-edition.
%   \editionsuffix  This edition's diagram-file suffix, e.g.
%                   \def\editionsuffix{.pl}. Defaults to empty. \dgm reads it,
%                   so a figure is called by slug and resolves per edition --
%                   see "Figures" below.
%
% A minimal English document:
%
%   \documentclass[11pt,a4paper]{article}
%   \newcommand{\paperstatus}{DRAFT}
%   \newcommand{\headerleft}{Project --- Overview}
%   \newcommand{\pdftitleline}{Project: an overview}
%   \newcommand{\pdfauthorline}{Author Name}
%   % ============================================================================
% HOUSE-PREAMBLE.tex -- the estate's house LaTeX style, as a file you copy in
% and \input rather than a block you copy out.
%
% See 00-RESEARCH-DOCUMENTATION.md, "The house preamble is a file, not a
% block". Copy this file to <repo>/docs/papers/house-preamble.tex (or
% docs/research/papers/ for a study paper) and \input it from every document
% in the repository that wants the house look. One copy per repository; every
% document in that repository \inputs the same one.
%
% WHY A FILE. The estate's first attempt at this shipped a template whose
% preamble was meant to be copied into each new document. Copying is what
% drifts: `ab-ove`'s two overview editions carried none of the seven house
% markers -- not the colors, the section formatting, the header, the status
% marker, the path macro, the geometry or the link style -- and nothing said
% so, because a convention about copying correctly has nothing to fail. The
% Beamer theme, which was always distributed as a file, propagated
% byte-identically over the same period. Distribute the file.
%
% ----------------------------------------------------------------------------
% THE CONTRACT. The including document defines these BEFORE \input, and
% carries its own \documentclass and \begin{document} -- this file has
% neither, so that one preamble can serve documents with different classes.
%
%   \paperstatus    Right-hand header, and the title block's status line.
%                   A study paper: DRAFT or VERIFIED, mirroring the study's
%                   own header. Any other document: its own status, named in
%                   that document's header comment as repurposed.
%   \headerleft     Left-hand header. Usually "Project --- Document title".
%   \pdftitleline   PDF metadata title.
%   \pdfauthorline  PDF metadata author.
%
% Optional, both defaulted below if the document does not define them:
%
%   \houselang      babel language for this edition, e.g. \def\houselang{polish}.
%                   Defaults to english. Set it in the edition file, not here:
%                   it is the one thing that is genuinely per-edition.
%   \editionsuffix  This edition's diagram-file suffix, e.g.
%                   \def\editionsuffix{.pl}. Defaults to empty. \dgm reads it,
%                   so a figure is called by slug and resolves per edition --
%                   see "Figures" below.
%
% A minimal English document:
%
%   \documentclass[11pt,a4paper]{article}
%   \newcommand{\paperstatus}{DRAFT}
%   \newcommand{\headerleft}{Project --- Overview}
%   \newcommand{\pdftitleline}{Project: an overview}
%   \newcommand{\pdfauthorline}{Author Name}
%   \input{house-preamble}
%   \begin{document}
%   ...
%
% Its Polish sibling differs by two lines, both before the \input:
%
%   \def\houselang{polish}
%   \def\editionsuffix{.pl}
%
% ----------------------------------------------------------------------------
% PACKAGES THIS FILE DOES NOT LOAD. The house set below is what every document
% needs. A document that needs more -- tikz, pifont, longtable, fancyvrb,
% amssymb -- loads it in its own preamble AFTER the \input, and says in a
% comment that it is an addition rather than part of the house set, so the
% next reader can tell the two apart at a glance.
% ============================================================================

% ----------------------------------------------------------------------------
% Encoding, then language. fontenc before inputenc before babel: babel reads
% the active font encoding when it sets up its shorthands, so loading it first
% gives an edition that is translated but typeset wrong.
% ----------------------------------------------------------------------------
\usepackage[T1]{fontenc}
\usepackage[utf8]{inputenc}

% The edition's language, defaulted so a monolingual document defines nothing.
% \PassOptionsToPackage rather than \usepackage[\houselang]{babel}: the option
% list here is a macro, and PassOptionsToPackage is the mechanism built for
% passing one.
\providecommand{\houselang}{english}
\PassOptionsToPackage{\houselang}{babel}
\usepackage{babel}

% ----------------------------------------------------------------------------
% The house package set.
% ----------------------------------------------------------------------------
\usepackage{geometry}
\usepackage{xcolor}
\usepackage{hyperref}
\usepackage{graphicx}
\usepackage{enumitem}
\usepackage{fancyhdr}
\usepackage{titlesec}
\usepackage{parskip}
\usepackage{booktabs}
\usepackage{array}
\usepackage{tabularx}

% ----------------------------------------------------------------------------
% Page geometry.
% ----------------------------------------------------------------------------
\geometry{
    a4paper,
    left=2.2cm,
    right=2.2cm,
    top=2.5cm,
    bottom=2.5cm
}

% ----------------------------------------------------------------------------
% House colors. Extracted from the estate's first formal documents
% (business_analysis.tex, pitch-deck-demium.tex) and unchanged since: a
% document from any repository should look like it came from the same shop.
% ----------------------------------------------------------------------------
\definecolor{primaryblue}{RGB}{41,128,185}
\definecolor{darkblue}{RGB}{31,78,121}
\definecolor{lightgray}{RGB}{236,240,241}
\definecolor{darkgray}{RGB}{52,73,94}
\definecolor{accentgreen}{RGB}{39,174,96}
\definecolor{accentorange}{RGB}{230,126,34}
\definecolor{accentred}{RGB}{231,76,60}

% ----------------------------------------------------------------------------
% Section formatting.
% ----------------------------------------------------------------------------
\titleformat{\section}
    {\color{primaryblue}\Large\bfseries\sffamily}
    {\thesection.}{0.5em}{}[\titlerule]

\titleformat{\subsection}
    {\color{darkblue}\large\bfseries\sffamily}
    {\thesubsection}{0.5em}{}

\titleformat{\subsubsection}
    {\color{darkgray}\normalsize\bfseries\sffamily}
    {\thesubsubsection}{0.5em}{}

% ----------------------------------------------------------------------------
% Header and footer. \paperstatus rides in the right-hand header of every
% page, so a draft cannot be mistaken for a final document at any page a
% reader happens to open it at -- which is the whole reason the marker is in
% the header rather than only on the title page.
% ----------------------------------------------------------------------------
\pagestyle{fancy}
\fancyhf{}
\renewcommand{\headrulewidth}{0.4pt}
\fancyhead[L]{\small\sffamily\color{darkgray} \headerleft}
\fancyhead[R]{\small\sffamily\color{accentred} \paperstatus\ --- \today}
\fancyfoot[C]{\thepage}

% ----------------------------------------------------------------------------
% Links and PDF metadata.
% ----------------------------------------------------------------------------
\hypersetup{
    colorlinks=true,
    linkcolor=primaryblue,
    urlcolor=primaryblue,
    pdfauthor={\pdfauthorline},
    pdftitle={\pdftitleline}
}

% ----------------------------------------------------------------------------
% Figures. A diagram's source is one .mmd file rendered to vector PDF by the
% repository's render step; the document includes the rendered file.
%
% \dgm resolves the path from a slug and the edition, so the body text names a
% diagram the way the diagrams directory does and never writes a path:
%
%   \includegraphics[width=0.78\linewidth]{\dgm{a1-system-context}}
%
% resolves to ../diagrams/rendered/a1-system-context.pdf in the default
% edition and ../diagrams/rendered/a1-system-context.pl.pdf in an edition that
% set \def\editionsuffix{.pl}. A repository whose editions share one set of
% diagrams simply never sets \editionsuffix, and every edition resolves to the
% same file.
%
% Hardcoding the path instead is what the estate did first, and it costs one
% edit per figure per edition: `ab-ove` maintained `b1-reader-loop.pdf` in the
% English file against `b1-reader-loop.pl.pdf` in the Polish one, by hand, in
% both. At three figures that is survivable; `agent-eval-bench` has 23.
%
% \housediagramdir is overridable for a document that does not sit one level
% below the diagrams directory -- define it before the \input.
% ----------------------------------------------------------------------------
\providecommand{\editionsuffix}{}
\providecommand{\housediagramdir}{../diagrams/rendered}
\newcommand{\dgm}[1]{\housediagramdir/#1\editionsuffix.pdf}

% ----------------------------------------------------------------------------
% House convenience commands.
%
% \repopath is for a path, a filename or an identifier copied from the code;
% it is deliberately the only one of these, so that a document does not grow a
% private synonym for it. \finding marks the sentence a section exists to
% deliver.
% ----------------------------------------------------------------------------
\newcommand{\repopath}[1]{\colorbox{lightgray}{\texttt{#1}}}
\newcommand{\finding}[1]{{\color{primaryblue}\textbf{#1}}}

% ----------------------------------------------------------------------------
% The house title block, as a command so that every document's first page
% looks the same and none of them reimplements it.
%
%   \housetitle{Title}{Subtitle or empty}{Author}{Affiliation or repo URL}{Provenance line}
%
% The provenance line is the one this estate keeps paying for the absence of:
% a document that presents another document names it here, so a reader who
% finds the two disagreeing knows which one is the source. For a study paper
% it names the companion study and the commit; for a project overview, the
% markdown it presents.
% ----------------------------------------------------------------------------
\newcommand{\housetitle}[5]{%
    \thispagestyle{plain}%
    \begin{center}
        {\LARGE\bfseries\sffamily\color{darkblue} #1}

        \ifx\relax#2\relax\else
            \vspace{0.3cm}

            {\Large\sffamily\color{darkgray} #2}
        \fi

        \vspace{0.6cm}

        {\large #3}

        \vspace{0.2cm}

        {\normalsize\color{darkgray} #4}

        \vspace{0.2cm}

        {\normalsize\color{accentred}\sffamily \paperstatus\ --- \today\ --- #5}
    \end{center}

    \vspace{0.4cm}
}

%   \begin{document}
%   ...
%
% Its Polish sibling differs by two lines, both before the \input:
%
%   \def\houselang{polish}
%   \def\editionsuffix{.pl}
%
% ----------------------------------------------------------------------------
% PACKAGES THIS FILE DOES NOT LOAD. The house set below is what every document
% needs. A document that needs more -- tikz, pifont, longtable, fancyvrb,
% amssymb -- loads it in its own preamble AFTER the \input, and says in a
% comment that it is an addition rather than part of the house set, so the
% next reader can tell the two apart at a glance.
% ============================================================================

% ----------------------------------------------------------------------------
% Encoding, then language. fontenc before inputenc before babel: babel reads
% the active font encoding when it sets up its shorthands, so loading it first
% gives an edition that is translated but typeset wrong.
% ----------------------------------------------------------------------------
\usepackage[T1]{fontenc}
\usepackage[utf8]{inputenc}

% The edition's language, defaulted so a monolingual document defines nothing.
% \PassOptionsToPackage rather than \usepackage[\houselang]{babel}: the option
% list here is a macro, and PassOptionsToPackage is the mechanism built for
% passing one.
\providecommand{\houselang}{english}
\PassOptionsToPackage{\houselang}{babel}
\usepackage{babel}

% ----------------------------------------------------------------------------
% The house package set.
% ----------------------------------------------------------------------------
\usepackage{geometry}
\usepackage{xcolor}
\usepackage{hyperref}
\usepackage{graphicx}
\usepackage{enumitem}
\usepackage{fancyhdr}
\usepackage{titlesec}
\usepackage{parskip}
\usepackage{booktabs}
\usepackage{array}
\usepackage{tabularx}

% ----------------------------------------------------------------------------
% Page geometry.
% ----------------------------------------------------------------------------
\geometry{
    a4paper,
    left=2.2cm,
    right=2.2cm,
    top=2.5cm,
    bottom=2.5cm
}

% ----------------------------------------------------------------------------
% House colors. Extracted from the estate's first formal documents
% (business_analysis.tex, pitch-deck-demium.tex) and unchanged since: a
% document from any repository should look like it came from the same shop.
% ----------------------------------------------------------------------------
\definecolor{primaryblue}{RGB}{41,128,185}
\definecolor{darkblue}{RGB}{31,78,121}
\definecolor{lightgray}{RGB}{236,240,241}
\definecolor{darkgray}{RGB}{52,73,94}
\definecolor{accentgreen}{RGB}{39,174,96}
\definecolor{accentorange}{RGB}{230,126,34}
\definecolor{accentred}{RGB}{231,76,60}

% ----------------------------------------------------------------------------
% Section formatting.
% ----------------------------------------------------------------------------
\titleformat{\section}
    {\color{primaryblue}\Large\bfseries\sffamily}
    {\thesection.}{0.5em}{}[\titlerule]

\titleformat{\subsection}
    {\color{darkblue}\large\bfseries\sffamily}
    {\thesubsection}{0.5em}{}

\titleformat{\subsubsection}
    {\color{darkgray}\normalsize\bfseries\sffamily}
    {\thesubsubsection}{0.5em}{}

% ----------------------------------------------------------------------------
% Header and footer. \paperstatus rides in the right-hand header of every
% page, so a draft cannot be mistaken for a final document at any page a
% reader happens to open it at -- which is the whole reason the marker is in
% the header rather than only on the title page.
% ----------------------------------------------------------------------------
\pagestyle{fancy}
\fancyhf{}
\renewcommand{\headrulewidth}{0.4pt}
\fancyhead[L]{\small\sffamily\color{darkgray} \headerleft}
\fancyhead[R]{\small\sffamily\color{accentred} \paperstatus\ --- \today}
\fancyfoot[C]{\thepage}

% ----------------------------------------------------------------------------
% Links and PDF metadata.
% ----------------------------------------------------------------------------
\hypersetup{
    colorlinks=true,
    linkcolor=primaryblue,
    urlcolor=primaryblue,
    pdfauthor={\pdfauthorline},
    pdftitle={\pdftitleline}
}

% ----------------------------------------------------------------------------
% Figures. A diagram's source is one .mmd file rendered to vector PDF by the
% repository's render step; the document includes the rendered file.
%
% \dgm resolves the path from a slug and the edition, so the body text names a
% diagram the way the diagrams directory does and never writes a path:
%
%   \includegraphics[width=0.78\linewidth]{\dgm{a1-system-context}}
%
% resolves to ../diagrams/rendered/a1-system-context.pdf in the default
% edition and ../diagrams/rendered/a1-system-context.pl.pdf in an edition that
% set \def\editionsuffix{.pl}. A repository whose editions share one set of
% diagrams simply never sets \editionsuffix, and every edition resolves to the
% same file.
%
% Hardcoding the path instead is what the estate did first, and it costs one
% edit per figure per edition: `ab-ove` maintained `b1-reader-loop.pdf` in the
% English file against `b1-reader-loop.pl.pdf` in the Polish one, by hand, in
% both. At three figures that is survivable; `agent-eval-bench` has 23.
%
% \housediagramdir is overridable for a document that does not sit one level
% below the diagrams directory -- define it before the \input.
% ----------------------------------------------------------------------------
\providecommand{\editionsuffix}{}
\providecommand{\housediagramdir}{../diagrams/rendered}
\newcommand{\dgm}[1]{\housediagramdir/#1\editionsuffix.pdf}

% ----------------------------------------------------------------------------
% House convenience commands.
%
% \repopath is for a path, a filename or an identifier copied from the code;
% it is deliberately the only one of these, so that a document does not grow a
% private synonym for it. \finding marks the sentence a section exists to
% deliver.
% ----------------------------------------------------------------------------
\newcommand{\repopath}[1]{\colorbox{lightgray}{\texttt{#1}}}
\newcommand{\finding}[1]{{\color{primaryblue}\textbf{#1}}}

% ----------------------------------------------------------------------------
% The house title block, as a command so that every document's first page
% looks the same and none of them reimplements it.
%
%   \housetitle{Title}{Subtitle or empty}{Author}{Affiliation or repo URL}{Provenance line}
%
% The provenance line is the one this estate keeps paying for the absence of:
% a document that presents another document names it here, so a reader who
% finds the two disagreeing knows which one is the source. For a study paper
% it names the companion study and the commit; for a project overview, the
% markdown it presents.
% ----------------------------------------------------------------------------
\newcommand{\housetitle}[5]{%
    \thispagestyle{plain}%
    \begin{center}
        {\LARGE\bfseries\sffamily\color{darkblue} #1}

        \ifx\relax#2\relax\else
            \vspace{0.3cm}

            {\Large\sffamily\color{darkgray} #2}
        \fi

        \vspace{0.6cm}

        {\large #3}

        \vspace{0.2cm}

        {\normalsize\color{darkgray} #4}

        \vspace{0.2cm}

        {\normalsize\color{accentred}\sffamily \paperstatus\ --- \today\ --- #5}
    \end{center}

    \vspace{0.4cm}
}

%   \begin{document}
%   ...
%
% Its Polish sibling differs by two lines, both before the \input:
%
%   \def\houselang{polish}
%   \def\editionsuffix{.pl}
%
% ----------------------------------------------------------------------------
% PACKAGES THIS FILE DOES NOT LOAD. The house set below is what every document
% needs. A document that needs more -- tikz, pifont, longtable, fancyvrb,
% amssymb -- loads it in its own preamble AFTER the \input, and says in a
% comment that it is an addition rather than part of the house set, so the
% next reader can tell the two apart at a glance.
% ============================================================================

% ----------------------------------------------------------------------------
% Encoding, then language. fontenc before inputenc before babel: babel reads
% the active font encoding when it sets up its shorthands, so loading it first
% gives an edition that is translated but typeset wrong.
% ----------------------------------------------------------------------------
\usepackage[T1]{fontenc}
\usepackage[utf8]{inputenc}

% The edition's language, defaulted so a monolingual document defines nothing.
% \PassOptionsToPackage rather than \usepackage[\houselang]{babel}: the option
% list here is a macro, and PassOptionsToPackage is the mechanism built for
% passing one.
\providecommand{\houselang}{english}
\PassOptionsToPackage{\houselang}{babel}
\usepackage{babel}

% ----------------------------------------------------------------------------
% The house package set.
% ----------------------------------------------------------------------------
\usepackage{geometry}
\usepackage{xcolor}
\usepackage{hyperref}
\usepackage{graphicx}
\usepackage{enumitem}
\usepackage{fancyhdr}
\usepackage{titlesec}
\usepackage{parskip}
\usepackage{booktabs}
\usepackage{array}
\usepackage{tabularx}

% ----------------------------------------------------------------------------
% Page geometry.
% ----------------------------------------------------------------------------
\geometry{
    a4paper,
    left=2.2cm,
    right=2.2cm,
    top=2.5cm,
    bottom=2.5cm
}

% ----------------------------------------------------------------------------
% House colors. Extracted from the estate's first formal documents
% (business_analysis.tex, pitch-deck-demium.tex) and unchanged since: a
% document from any repository should look like it came from the same shop.
% ----------------------------------------------------------------------------
\definecolor{primaryblue}{RGB}{41,128,185}
\definecolor{darkblue}{RGB}{31,78,121}
\definecolor{lightgray}{RGB}{236,240,241}
\definecolor{darkgray}{RGB}{52,73,94}
\definecolor{accentgreen}{RGB}{39,174,96}
\definecolor{accentorange}{RGB}{230,126,34}
\definecolor{accentred}{RGB}{231,76,60}

% ----------------------------------------------------------------------------
% Section formatting.
% ----------------------------------------------------------------------------
\titleformat{\section}
    {\color{primaryblue}\Large\bfseries\sffamily}
    {\thesection.}{0.5em}{}[\titlerule]

\titleformat{\subsection}
    {\color{darkblue}\large\bfseries\sffamily}
    {\thesubsection}{0.5em}{}

\titleformat{\subsubsection}
    {\color{darkgray}\normalsize\bfseries\sffamily}
    {\thesubsubsection}{0.5em}{}

% ----------------------------------------------------------------------------
% Header and footer. \paperstatus rides in the right-hand header of every
% page, so a draft cannot be mistaken for a final document at any page a
% reader happens to open it at -- which is the whole reason the marker is in
% the header rather than only on the title page.
% ----------------------------------------------------------------------------
\pagestyle{fancy}
\fancyhf{}
\renewcommand{\headrulewidth}{0.4pt}
\fancyhead[L]{\small\sffamily\color{darkgray} \headerleft}
\fancyhead[R]{\small\sffamily\color{accentred} \paperstatus\ --- \today}
\fancyfoot[C]{\thepage}

% ----------------------------------------------------------------------------
% Links and PDF metadata.
% ----------------------------------------------------------------------------
\hypersetup{
    colorlinks=true,
    linkcolor=primaryblue,
    urlcolor=primaryblue,
    pdfauthor={\pdfauthorline},
    pdftitle={\pdftitleline}
}

% ----------------------------------------------------------------------------
% Figures. A diagram's source is one .mmd file rendered to vector PDF by the
% repository's render step; the document includes the rendered file.
%
% \dgm resolves the path from a slug and the edition, so the body text names a
% diagram the way the diagrams directory does and never writes a path:
%
%   \includegraphics[width=0.78\linewidth]{\dgm{a1-system-context}}
%
% resolves to ../diagrams/rendered/a1-system-context.pdf in the default
% edition and ../diagrams/rendered/a1-system-context.pl.pdf in an edition that
% set \def\editionsuffix{.pl}. A repository whose editions share one set of
% diagrams simply never sets \editionsuffix, and every edition resolves to the
% same file.
%
% Hardcoding the path instead is what the estate did first, and it costs one
% edit per figure per edition: `ab-ove` maintained `b1-reader-loop.pdf` in the
% English file against `b1-reader-loop.pl.pdf` in the Polish one, by hand, in
% both. At three figures that is survivable; `agent-eval-bench` has 23.
%
% \housediagramdir is overridable for a document that does not sit one level
% below the diagrams directory -- define it before the \input.
% ----------------------------------------------------------------------------
\providecommand{\editionsuffix}{}
\providecommand{\housediagramdir}{../diagrams/rendered}
\newcommand{\dgm}[1]{\housediagramdir/#1\editionsuffix.pdf}

% ----------------------------------------------------------------------------
% House convenience commands.
%
% \repopath is for a path, a filename or an identifier copied from the code;
% it is deliberately the only one of these, so that a document does not grow a
% private synonym for it. \finding marks the sentence a section exists to
% deliver.
% ----------------------------------------------------------------------------
\newcommand{\repopath}[1]{\colorbox{lightgray}{\texttt{#1}}}
\newcommand{\finding}[1]{{\color{primaryblue}\textbf{#1}}}

% ----------------------------------------------------------------------------
% The house title block, as a command so that every document's first page
% looks the same and none of them reimplements it.
%
%   \housetitle{Title}{Subtitle or empty}{Author}{Affiliation or repo URL}{Provenance line}
%
% The provenance line is the one this estate keeps paying for the absence of:
% a document that presents another document names it here, so a reader who
% finds the two disagreeing knows which one is the source. For a study paper
% it names the companion study and the commit; for a project overview, the
% markdown it presents.
% ----------------------------------------------------------------------------
\newcommand{\housetitle}[5]{%
    \thispagestyle{plain}%
    \begin{center}
        {\LARGE\bfseries\sffamily\color{darkblue} #1}

        \ifx\relax#2\relax\else
            \vspace{0.3cm}

            {\Large\sffamily\color{darkgray} #2}
        \fi

        \vspace{0.6cm}

        {\large #3}

        \vspace{0.2cm}

        {\normalsize\color{darkgray} #4}

        \vspace{0.2cm}

        {\normalsize\color{accentred}\sffamily \paperstatus\ --- \today\ --- #5}
    \end{center}

    \vspace{0.4cm}
}


\setlength{\parindent}{0pt}
\setlist[enumerate,1]{label=\arabic*.,leftmargin=*,itemsep=2pt}

\begin{document}

\housetitle{Pilotaż 12 miesięcy: cennik i wzór umowy}
  {Prywatna sztuczna inteligencja w gabinecie, urządzenie offline}
  {Konrad Cinkusz, Inference Polska}
  {[adres], NIP [numer]}
  {cennik pilotażowy i wzór umowy najmu}

\section*{Cennik pilotażowy}

Pilotaż wychodzi na zero przy opłacie instalacyjnej 1\,500 zł netto i abonamencie 590 zł netto miesięcznie przez 12 miesięcy: razem 8\,580 zł netto przy kosztach ok. 8\,390 zł.

\subsection*{Koszty 12 miesięcy (netto)}
\begin{tabularx}{\linewidth}{@{}Xrl@{}}
\toprule
Pozycja & Kwota & Podstawa \\
\midrule
Utrata wartości sprzętu & 4\,990 zł & zakup ok. 9\,990 zł netto, odsprzedaż lub wykup po roku za 50\% \\
Rezerwa na serwis i wymianę & 1\,000 zł & 10\% wartości sprzętu rocznie \\
Czas Wynajmującego & 2\,400 zł & 12 h wsparcia po 200 zł \\
\midrule
\textbf{Suma kosztów} & \textbf{8\,390 zł} & \\
\bottomrule
\end{tabularx}

\subsection*{Trzy warianty opłat, każdy na zero (netto, przed VAT 23\%)}
\begin{tabularx}{\linewidth}{@{}Xrrrrr@{}}
\toprule
Wariant & Instalacja & Abonament $\times$ 12 & Razem & Brutto: instalacja & Brutto: abonament \\
\midrule
A: abonament 500 zł brutto & 3\,500 zł & 407 zł & 8\,384 zł & 4\,305 zł & 500 zł \\
B: zrównoważony (polecany) & 1\,500 zł & 590 zł & 8\,580 zł & 1\,845 zł & 726 zł \\
C: bez opłaty instalacyjnej & 0 zł & 700 zł & 8\,400 zł & 0 zł & 861 zł \\
\bottomrule
\end{tabularx}

Wariant B rozkłada koszt: umiarkowana opłata na starcie i abonament poniżej 750 zł brutto. Wariant A utrzymuje 500 zł brutto miesięcznie, ale klient płaci 4\,305 zł brutto na początku. Wariant C nie ma opłaty wstępnej, za to abonament jest najwyższy.

\subsection*{Poza abonamentem (netto)}
\begin{tabularx}{\linewidth}{@{}Xr@{}}
\toprule
Usługa & Cena \\
\midrule
Dodatkowa godzina wsparcia & 200 zł \\
Wizyta serwisowa na miejscu & 250 zł plus dojazd \\
Dodatkowe szkolenie zespołu (2 h) & 400 zł \\
Wykup urządzenia po 12 miesiącach & 5\,000 zł \\
Opłata za wcześniejsze zakończenie & 50\% pozostałych abonamentów \\
\bottomrule
\end{tabularx}

\paragraph{Założenia.} Klient płaci za prąd. Gabinet leży w zasięgu bez kosztu dojazdu. Jesteś płatnikiem VAT. Wartość sprzętu po roku to 50\% ceny zakupu (nie zweryfikowane). 12 h pracy to 1 h wsparcia miesięcznie. Przy droższym RAM-ie (Kingston, 14\,683 zł brutto zamiast 12\,284 zł) koszty rosną do ok. 9\,560 zł, a abonament w wariancie B do ok. 670 zł netto.

\paragraph{Ceny sprzętu z 14.11.2026:} RTX 5070, x-kom, RAM 64 GB Crucial, Morele, RAM 64 GB Kingston, x-kom.

\clearpage
\section*{Umowa najmu urządzenia i usług wsparcia (pilotaż)}

zawarta w dniu [data] w [miejscowość] pomiędzy:

Konradem Cinkuszem, prowadzącym działalność gospodarczą pod firmą [pełna nazwa firmy z CEIDG], NIP [NIP], REGON [REGON], adres [adres], e-mail [e-mail], zwanym dalej Wynajmującym,

a

[nazwa gabinetu lub podmiotu], NIP [NIP], adres [adres], reprezentowanym przez [imię i nazwisko, funkcja], zwanym dalej Najemcą.

\subsection*{§1 Przedmiot umowy}
\begin{enumerate}
\item Wynajmujący oddaje Najemcy do używania urządzenie opisane w Załączniku nr 1 (dalej: Urządzenie) wraz z zainstalowanym oprogramowaniem i lokalnymi modelami językowymi (dalej: Oprogramowanie) oraz świadczy usługi wsparcia opisane w §4.
\item Urządzenie jest przeznaczone do pracy bez dostępu do sieci (offline) i służy do wspomagania pracy Najemcy przy analizie i redagowaniu dokumentów.
\item Umowa ma charakter pilotażowy: poza najmem służy zebraniu doświadczeń do opisu studium przypadku zgodnie z §12.
\end{enumerate}

\subsection*{§2 Czas trwania}
\begin{enumerate}
\item Umowa zostaje zawarta na czas określony 12 miesięcy, liczony od dnia podpisania protokołu odbioru Urządzenia (Załącznik nr 3), zwany dalej Okresem Pilotażu.
\item Po zakończeniu Okresu Pilotażu umowa wygasa. Nie później niż 30 dni przed jego końcem Strony mogą uzgodnić na piśmie przedłużenie najmu albo wykup Urządzenia zgodnie z §15.
\item Bez takiego uzgodnienia Najemca zwraca Urządzenie zgodnie z §14.
\end{enumerate}

\subsection*{§3 Instalacja i odbiór}
\begin{enumerate}
\item Wynajmujący dostarcza, instaluje i konfiguruje Urządzenie w siedzibie Najemcy w terminie [termin] od zawarcia umowy.
\item Odbiór potwierdza protokół podpisany przez obie Strony (Załącznik nr 3). Od jego podpisania liczy się Okres Pilotażu i abonament.
\item Wynajmujący przeprowadza szkolenie zespołu Najemcy w wymiarze 2 godzin, w terminie ustalonym z Najemcą.
\item Najemca wskazuje miejsce ustawienia Urządzenia spełniające wymagania z §9.
\end{enumerate}

\subsection*{§4 Zakres usług wsparcia}
\begin{enumerate}
\item W ramach abonamentu Wynajmujący zapewnia: aktualizacje Oprogramowania i baz wiedzy, serwis i wymianę uszkodzonego sprzętu oraz do 1 godziny wsparcia miesięcznie.
\item Urządzenie działa offline, więc aktualizacje wykonuje Wynajmujący na miejscu, z zweryfikowanych nośników, w terminie uzgodnionym z Najemcą.
\item Wsparcie ponad limit z ust. 1, wizyty serwisowe oraz dodatkowe szkolenia są płatne według Załącznika nr 1.
\item Awarię zgłasza się na adres [e-mail] lub pod numer [telefon]. Wynajmujący podejmuje naprawy lub wymienia sprzęt w ciągu [5] dni roboczych od zgłoszenia.
\end{enumerate}

\subsection*{§5 Parametry i ograniczenia działania}
\begin{enumerate}
\item Urządzenie obsługuje jednego użytkownika na raz. Czas odpowiedzi modeli zależy od modelu i długości dokumentu; wartości orientacyjne podaje Załącznik nr 1.
\item Wynajmujący nie gwarantuje określonej dostępności ani czasu odpowiedzi. Umowa nie zawiera umowy o poziomie usług (SLA).
\end{enumerate}

\subsection*{§6 Tryb offline i bezpieczeństwo}
\begin{enumerate}
\item Urządzenie pracuje bez połączenia z siecią. Najemca nie podłącza go do sieci przewodowej ani bezprzewodowej, do internetu ani do innych urządzeń poza drukarką i skanerem wskazanymi w Załączniku nr 1.
\item Najemca nie wprowadza do Urządzenia nośników danych niezweryfikowanych lub nieznanego pochodzenia.
\item Aktualizacje Oprogramowania i baz wykonuje wyłącznie Wynajmujący, osobiście na miejscu, z zweryfikowanych nośników. Wynajmujący nie ma zdalnego dostępu do Urządzenia.
\item Najemca nie modyfikuje Oprogramowania, nie instaluje innego oprogramowania i nie otwiera obudowy Urządzenia.
\item Dostęp do Urządzenia mają wyłącznie osoby upoważnione przez Najemcę. Najemca odpowiada za konta i hasła oraz za zabezpieczenie Urządzenia przed dostępem osób trzecich (zamknięte pomieszczenie).
\item Wynajmujący konfiguruje szyfrowanie dysku z danymi.
\end{enumerate}

\subsection*{§7 Dane osobowe i tajemnica zawodowa}
\begin{enumerate}
\item Dane przetwarzane na Urządzeniu, w tym dane o stanie zdrowia pacjentów, pozostają w wyłącznej dyspozycji Najemcy, który jest ich administratorem w rozumieniu RODO. Wynajmujący nie kopiuje ich, nie wynosi i nie przetwarza we własnych celach.
\item W zakresie, w jakim podczas instalacji, serwisu lub aktualizacji Wynajmujący może uzyskać dostęp do danych osobowych, Strony zawierają umowę powierzenia przetwarzania danych zgodną z art. 28 RODO (Załącznik nr 2).
\item Uszkodzony nośnik danych nie opuszcza siedziby Najemcy bez jego pisemnej zgody. Po wymianie nośnik jest przekazywany Najemcy albo trwale niszczony w jego obecności.
\item Najemca odpowiada za zgodność własnego korzystania z Urządzenia z przepisami, w szczególności za podstawę prawną przetwarzania danych, tajemnicę zawodową i informowanie pacjentów, jeśli jest to wymagane.
\item Po zakończeniu umowy dane usuwa się zgodnie z §14.
\end{enumerate}

\subsection*{§8 Wynagrodzenie}
\begin{enumerate}
\item Najemca płaci Wynajmującemu: (a) jednorazową opłatę instalacyjną w wysokości [1\,500] zł netto, (b) abonament miesięczny w wysokości [590] zł netto za każdy miesiąc Okresu Pilotażu. Kwoty zwiększa się o podatek VAT według stawki obowiązującej w dniu wystawienia faktury.
\item Fakturę na opłatę instalacyjną Wynajmujący wystawia w dniu podpisania protokołu odbioru; termin płatności to 14 dni. Faktury na abonament wystawia się z dołu do 5. dnia miesiąca; termin płatności to 14 dni. Za pierwszy, niepełny miesiąc abonament nalicza się proporcjonalnie do liczby dni.
\item Faktury wystawia się zgodnie z obowiązującymi przepisami, w tym przez KSeF, jeśli ma on zastosowanie.
\item Za opóźnienie w płatności Wynajmujący nalicza odsetki ustawowe za opóźnienie w transakcjach handlowych. Po 14 dniach zwłoki i bezskutecznym wezwaniu Wynajmujący może wstrzymać świadczenie usług wsparcia.
\item W Okresie Pilotażu ceny nie podlegają waloryzacji.
\item Usługi dodatkowe z §4 ust. 3 są płatne odrębnie, według Załącznika nr 1.
\end{enumerate}

\subsection*{§9 Obowiązki Najemcy}
\begin{enumerate}
\item Najemca zapewnia miejsce i zasilanie Urządzenia (230 V), odpowiednią wentylację i temperaturę, ochronę przed zalaniem i kradzieżą oraz ponosi koszt energii elektrycznej.
\item Najemca używa Urządzenia zgodnie z przeznaczeniem i instrukcją, nie podnajmuje go ani nie użycza osobom trzecim.
\item Najemca niezwłocznie, nie później niż w ciągu 24 godzin od wykrycia, zgłasza awarię, uszkodzenie, utratę Urządzenia, nieupoważniony dostęp lub podejrzenie incydentu bezpieczeństwa.
\item Najemca zapewnia Wynajmującemu dostęp do Urządzenia w terminach uzgodnionych na potrzeby serwisu i aktualizacji.
\item Osobą kontaktową Najemcy jest [imię i nazwisko, telefon, e-mail].
\end{enumerate}

\subsection*{§10 Własność, ryzyko i szkody}
\begin{enumerate}
\item Urządzenie i Oprogramowanie pozostają własnością Wynajmującego. Najemca nie może ich zbywać, obciążać ani udostępniać osobom trzecim. Oprogramowanie własne Wynajmującego jest licencjonowane niewyłącznie na czas umowy; modele językowe działają na zasadach swoich otwartych licencji.
\item Od podpisania protokołu odbioru do zwrotu ryzyko utraty, kradzieży lub zniszczenia Urządzenia w miejscu użytkowania, z przyczyn niezależnych od Wynajmującego, ponosi Najemca. [Najemca obejmuje Urządzenie ubezpieczeniem mienia od kradzieży i zalania na kwotę wskazaną w Załączniku nr 1.]
\item Najemca pokrywa szkodę ponad normalne zużycie wynikające z prawidłowego używania, do wysokości wartości Urządzenia wskazanej w Załączniku nr 1.
\end{enumerate}

\subsection*{§11 Wyniki modeli AI i odpowiedzialność}
\begin{enumerate}
\item Najemca przyjmuje do wiadomości, że modele językowe mogą generować treści nieprawdziwe, niepełne lub nieadekwatne. Wynik jest materiałem pomocniczym, który musi zweryfikować uprawniony specjalista po stronie Najemcy.
\item Urządzenie i usługi nie stanowią porady medycznej, prawnej ani podatkowej, nie są wyrobem medycznym i nie są przeznaczone do diagnozowania ani leczenia. Najemca nie opiera decyzji diagnostycznych ani terapeutycznych wyłącznie na wynikach modeli.
\item Odpowiedzialność Wynajmującego z tytułu umowy, z wyłączeniem szkody wyrządzonej umyślnie, ogranicza się do kwoty wynagrodzenia zapłaconego przez Najemcę w ciągu ostatnich 12 miesięcy. Wynajmujący nie odpowiada za utracone korzyści.
\item Wynajmujący nie odpowiada za szkody wynikające z użycia Urządzenia niezgodnie z umową lub z braku weryfikacji wyników przez Najemcę.
\end{enumerate}

\subsection*{§12 Studium przypadku}
\begin{enumerate}
\item Najemca zgadza się, aby Wynajmujący opisał przebieg pilotażu (cele, konfigurację, czas pracy, koszty i wnioski o użyteczności) w formie anonimowego studium przypadku, publikowanego na stronie Wynajmującego i w jego materiałach.
\item Studium nie zawiera danych osobowych pacjentów ani treści dokumentacji. Wynajmujący nie otrzymuje dokumentacji Najemcy do wglądu. Podstawą studium są informacje przekazywane dobrowolnie przez Najemcę (ankieta, liczba użyć, wrażenia).
\item Nazwa Najemcy i jego opinia o usłudze mogą zostać ujawnione tylko za odrębną zgodą Najemcy w formie dokumentowej, którą można w każdej chwili cofnąć na przyszłość.
\item Strony spotykają się po 3, 6 i 12 miesiącu Okresu Pilotażu na rozmowę podsumowującą (ok. 30 minut).
\end{enumerate}

\subsection*{§13 Rozwiązanie umowy}
\begin{enumerate}
\item Najemca może wypowiedzieć umowę z zachowaniem 30-dniowego okresu wypowiedzenia ze skutkiem na koniec miesiąca kalendarzowego. W takim przypadku płaci opłatę za wcześniejsze zakończenie w wysokości 50\% sumy abonamentów przypadających do końca Okresu Pilotażu. Opłata nie należy się, gdy przyczyną wypowiedzenia jest rażące naruszenie umowy przez Wynajmującego.
\item Wynajmujący może wypowiedzieć umowę z zachowaniem 30-dniowego okresu wypowiedzenia z ważnej przyczyny, w szczególności przy zwłoce w płatności dłuższej niż 30 dni po wezwaniu albo używaniu Urządzenia niezgodnie z przeznaczeniem.
\item Rażące naruszenie §6 ust. 1–4 uprawnia Wynajmującego do rozwiązania umowy bez zachowania okresu wypowiedzenia.
\item Wypowiedzenie wymaga formy pisemnej lub dokumentowej.
\end{enumerate}

\subsection*{§14 Zwrot Urządzenia i usunięcie danych}
\begin{enumerate}
\item Najemca wydaje Urządzenie Wynajmującemu w siedzibie Najemcy w ciągu 7 dni od zakończenia umowy.
\item Przed zwrotem Najemca może zabezpieczyć własne dane (wydruki, kopie na swoim nośniku). Na jego życzenie Wynajmujący pomaga w eksporcie.
\item Dane Najemcy są trwale usuwane z Urządzenia (bezpieczne wymazanie lub kasowanie kryptograficzne) w obecności przedstawiciela Najemcy, co potwierdza protokół z Załącznika nr 4. Wynajmujący nie wynosi żadnych danych Najemcy.
\item Za każdy dzień zwłoki ze zwrotem Najemca płaci wynagrodzenie za korzystanie z Urządzenia proporcjonalnie do abonamentu.
\end{enumerate}

\subsection*{§15 Wykup}
\begin{enumerate}
\item Najemca może, nie później niż 30 dni przed końcem Okresu Pilotażu, złożyć pisemne oświadczenie o wykupie Urządzenia. Wykup nie jest obowiązkiem.
\item Cena wykupu wynosi [5\,000] zł netto plus VAT (Załącznik nr 1). Własność sprzętu przechodzi na Najemcę z chwilą zapłaty.
\item Wykup obejmuje sprzęt i modele językowe na ich otwartych licencjach. Oprogramowanie własne Wynajmującego pozostaje jego własnością, chyba że Strony zawrą odrębną umowę licencyjną lub wsparcia.
\end{enumerate}

\subsection*{§16 Poufność}
\begin{enumerate}
\item Strony zachowują w poufności informacje techniczne i handlowe drugiej Strony, w tym warunki umowy.
\item Wynajmujący zachowuje w tajemnicy informacje o Najemcy, jego pacjentach i klientach, które mógłby poznać w związku z umową, również po jej zakończeniu, bez ograniczenia w czasie.
\item Zobowiązanie nie dotyczy ujawnienia wymaganego przepisami prawa ani publikacji, na którą Najemca wyraził zgodę na podstawie §12.
\end{enumerate}

\subsection*{§17 Postanowienia końcowe}
\begin{enumerate}
\item Zmiany umowy wymagają formy pisemnej pod rygorem nieważności.
\item Załączniki są integralną częścią umowy. W razie sprzeczności pierwszeństwo ma treść umowy, a w sprawach ochrony danych osobowych Załącznik nr 2.
\item Żadna ze Stron nie odpowiada za niewykonanie zobowiązania spowodowane siłą wyższą.
\item Do umowy stosuje się prawo polskie. Spory rozstrzyga sąd właściwy dla [siedziby Wynajmującego].
\item Umowę sporządzono w dwóch jednobrzmiących egzemplarzach, po jednym dla każdej ze Stron.
\end{enumerate}

\vspace{1.5cm}
\begin{tabularx}{\linewidth}{@{}XX@{}}
[miejsce, data] & [miejsce, data] \\[2.2cm]
\hrule & \hrule \\
Wynajmujący: Konrad Cinkusz & Najemca: [imię i nazwisko, funkcja] \\
\end{tabularx}

\clearpage
\section*{Załącznik nr 1: specyfikacja i cennik}

\begin{tabularx}{\linewidth}{@{}>{\bfseries}p{0.34\linewidth}X@{}}
\toprule
Parametr & Wartość \\
\midrule
Urządzenie & stacja robocza, nr seryjny [numer] \\
Karta graficzna & 12 GB VRAM, [model] \\
Procesor & [model] \\
Pamięć RAM & 64 GB DDR5 \\
Dysk & NVMe [1 TB], szyfrowany \\
System & Linux \\
Interfejsy sieciowe & wyłączone; Wi-Fi i Bluetooth niezamontowane \\
Peryferia & drukarka i skaner [model] \\
Modele językowe & Qwen3.5 9B, 27B i 35B-A3B (Apache 2.0 według producenta; sprawdź kartę modelu) \\
Obciążenie & jeden użytkownik na raz \\
Czas odpowiedzi (orientacyjnie) & [uzupełnić po pomiarach] \\
Wartość odtworzeniowa Urządzenia & [ok. 10\,000] zł netto \\
Ubezpieczenie & [wymagane / zalecane], suma [kwota] zł \\
\bottomrule
\end{tabularx}

\vspace{0.6cm}
\begin{tabularx}{\linewidth}{@{}Xrr@{}}
\toprule
Pozycja & Netto & Brutto (VAT 23\%) \\
\midrule
Opłata instalacyjna (jednorazowo) & 1\,500 zł & 1\,845 zł \\
Abonament miesięczny (12 miesięcy) & 590 zł & 726 zł \\
Wykup po Okresie Pilotażu & 5\,000 zł & 6\,150 zł \\
Dodatkowa godzina wsparcia & 200 zł & 246 zł \\
Wizyta serwisowa (plus dojazd [stawka]) & 250 zł & 308 zł \\
Dodatkowe szkolenie zespołu (2 h) & 400 zł & 492 zł \\
Opłata za wcześniejsze zakończenie (§13) & \multicolumn{2}{r}{50\% pozostałych abonamentów} \\
\bottomrule
\end{tabularx}

\clearpage
\section*{Załącznik nr 2: umowa powierzenia przetwarzania danych (art. 28 RODO), szkic}

\begin{enumerate}
\item Administratorem jest Najemca, podmiotem przetwarzającym Wynajmujący. Powierzenie obejmuje wyłącznie dostęp do danych, jaki może wystąpić podczas instalacji, serwisu i aktualizacji Urządzenia. Wynajmujący nie ma zdalnego dostępu do danych, a prace wykonuje w obecności przedstawiciela Najemcy.
\item Rodzaje danych: dane zwykłe i dane szczególnych kategorii (dane o zdrowiu). Kategorie osób: pacjenci, pracownicy i współpracownicy Najemcy.
\item Wynajmujący przetwarza dane wyłącznie na udokumentowane polecenie Najemcy; umowa wraz z załącznikami jest takim poleceniem.
\item Osoby upoważnione po stronie Wynajmującego zachowują poufność. Wynajmujący stosuje środki bezpieczeństwa z art. 32 RODO, w tym szyfrowanie dysku.
\item Podpowierzenie przetwarzania jest wyłączone, chyba że Najemca zgodzi się na nie pisemnie.
\item Wynajmujący pomaga Najemcy w realizacji praw osoby, której dane dotyczą, i w obowiązkach z art. 32–36 RODO.
\item Wynajmujący zgłasza Najemcy naruszenie ochrony danych bez zbędnej zwłoki, nie później niż w ciągu 24 godzin od stwierdzenia.
\item Po zakończeniu świadczenia usług Wynajmujący usuwa dane zgodnie z §14 umowy i potwierdza to protokołem.
\item Wynajmujący udostępnia Najemcy informacje potrzebne do wykazania zgodności z art. 28 RODO i umożliwia audyt w ustalonym terminie. Informuje o poleceniu, które uważa za naruszające przepisy.
\end{enumerate}

\section*{Załącznik nr 3: protokół odbioru Urządzenia}

\begin{tabularx}{\linewidth}{@{}>{\bfseries}p{0.34\linewidth}X@{}}
\toprule
Pole & Wpis \\
\midrule
Data i miejsce & [ ] \\
Numer seryjny Urządzenia & [ ] \\
Stan wizualny i kompletność & [ ] \\
Test działania modeli & [ ] \\
Interfejsy sieciowe wyłączone & [tak / nie] \\
Szkolenie 2 h przeprowadzone & [tak / nie] \\
Uwagi & [ ] \\
Podpisy & Wynajmujący: [ ] \quad Najemca: [ ] \\
\bottomrule
\end{tabularx}

\section*{Załącznik nr 4: protokół zwrotu i usunięcia danych}

\begin{tabularx}{\linewidth}{@{}>{\bfseries}p{0.34\linewidth}X@{}}
\toprule
Pole & Wpis \\
\midrule
Data i miejsce & [ ] \\
Numer seryjny Urządzenia & [ ] \\
Stan Urządzenia przy zwrocie & [ ] \\
Metoda usunięcia danych & [bezpieczne wymazanie / kasowanie kryptograficzne] \\
Usunięcie wykonano w obecności & [imię i nazwisko] \\
Dane Najemcy zabezpieczone przez Najemcę przed zwrotem & [tak / nie] \\
Podpisy & Wynajmujący: [ ] \quad Najemca: [ ] \\
\bottomrule
\end{tabularx}

\end{document}
